\documentclass[11pt]{beamer}
\usetheme{Madrid}

\usepackage{fontspec}
\usepackage[spanish,es-nodecimaldot]{babel}
\usepackage{amsmath}
\usepackage{amsfonts}
\usepackage{amssymb}
\usepackage{animate}
\usepackage{graphicx}
\usepackage{hyperref}
\usepackage{color}

\author{José Ignacio Escribano}
\title{Superficies regladas}
\date{23 de abril de 2015}
\institute[URJC]{\includegraphics[width=0.1\linewidth]{logoURJC.jpg}}

\setbeamertemplate{caption}[numbered]

\usepackage{amsmath}
\usepackage{xparse}
\usepackage{amsthm}
\usepackage{tikz}

% Seno
\def\sen{\mathop{\mbox{\normalfont sen}}\nolimits}

% Seno hiperbólico
\def\senh{\mathop{\mbox{\normalfont senh}}\nolimits}

% Arcoseno
\def\arcsen{\mathop{\mbox{\normalfont arcsen}}\nolimits}

% Gradiente
\def\grad{\mathop{\mbox{\normalfont grad}}\nolimits}

% Argumento principal
\def\Arg{\mathop{\mbox{\normalfont Arg}}\nolimits}

% Logaritmo complejo
\def\Log{\mathop{\mbox{\normalfont Log}}\nolimits}

% Rotacional
\def\rot{\mathop{\mbox{\normalfont rot}}\nolimits}

% Divergencia
\def\diver{\mathop{\mbox{\normalfont div}}\nolimits}

% Longitud
\def\longi{\mathop{\mbox{\normalfont long}}\nolimits}

% Índice de una curva
\def\ind{\mathop{\mbox{\normalfont Ind}}\nolimits}

% Norma
\newcommand{\norm}[1]{\lVert #1 \rVert}

% Derivada parcial
\newcommand{\derpar}[1]{\dfrac{\partial}{\partial #1}}

% Derivada parcial con función
\newcommand{\derparcial}[2]{\dfrac{\partial #1}{\partial #2}}

% Derivada parcial de orden n con funcion
\newcommand{\derparcialn}[3]{\dfrac{\partial^{#3} #1}{\partial #2^{#3}}}

% Integral doble en [a,b] x [c,d]
\newcommand{\intdob}[4]{\displaystyle \int_{#1}^{#2} \int_{#3}^{#4}}

% Integral triple en [a,b] x [c,d] x [e,f]
\newcommand{\inttri}[6]{\displaystyle \int_{#1}^{#2} \int_{#3}^{#4} \int_{#5}^{#6}}

% Números complejos
\def\C{\ensuremath{\mathbb{C}}}

% Números reales
\def\R{\ensuremath{\mathbb{R}}}

% Números racionales
\def\Q{\ensuremath{\mathbb{Q}}}

% Números enteros
\def\Z{\ensuremath{\mathbb{Z}}}

% Números naturales
\def\N{\ensuremath{\mathbb{N}}}

% Teorema
\newtheorem{teo}{Teorema}

% Corolario
\newtheorem{cor}{Corolario}

% Proposición
\newtheorem{prop}{Proposición}

% Lema
\newtheorem{lema}{Lema}

% Definición
\newtheorem{defi}{Definición}

% Observaciones
\newtheorem*{obs}{Observaciones}

% Observación
\newtheorem*{ob}{Observación}

% Nota
\newtheorem*{nota}{Nota}

% Pregunta
\newtheorem*{pregunta}{Pregunta}

% Ejemplo
\newtheorem*{ejemplo}{Ejemplo}


\NewDocumentCommand{\overarrow}{O{=} O{\uparrow} m}{%
  \overset{\makebox[0pt]{\begin{tabular}{@{}c@{}}#3\\[0pt]\ensuremath{#2}\end{tabular}}}{#1}
}
\NewDocumentCommand{\underarrow}{O{=} O{\downarrow} m}{%
  \underset{\makebox[0pt]{\begin{tabular}{@{}c@{}}\ensuremath{#2}\\[0pt]#3\end{tabular}}}{\ensuremath{#1}}
}

\newcommand{\contradiction}{%
\begin{tikzpicture}[rotate=45,x=0.5ex,y=0.5ex]
\draw[color=red, line width=.1ex] (0,2) -- (3,2) (0,1) -- (3,1) (1,3) -- (1,0) (2,3) -- (2,0);
\end{tikzpicture}
}


\begin{document}
	\begin{frame}[plain]
		\maketitle
	\end{frame}
	
	\begin{frame}{Índice}
		\tableofcontents
	\end{frame}
	
	
	\section{Introducción}
	
	\begin{frame}
		\begin{center}
			\Huge\textbf{\textsf{\textcolor{blue}{Introducción}}}
		\end{center}
	\end{frame}
	
	
	\begin{frame}{Introducción}
		\begin{defi}
			Una superficie es una aplicación
			\begin{equation}
			r : [a,b] \times [c,d] \subseteq \R^2 \to \R^3
			\end{equation}
		\end{defi}
		
		\begin{ejemplo}[de superficie]
			\textcolor{red}{Imagen de superficie}
		\end{ejemplo}	
	\end{frame}
	
	\begin{frame}{Introducción}
		\begin{defi}
			Una superficie reglada es aquella generada por una recta en el espacio llamada generatriz, a lo largo de una curva llamada directriz.
		\end{defi}	
		
		\begin{defi}
			Una superficie es reglada si es generada por una familia uniparamétrica de rectas (llamadas generatrices).\\
			
			La parametrización de una superficie reglada es 
			
			\begin{equation}
			r(u,v) = \rho(u) + v a(u)
			\end{equation}
			
			donde $\rho, a : I = [a,b] \to \R^3$ son dos curvas en el espacio.
		\end{defi}
		
		\begin{defi}
			Un conjunto de curvas de una superficie reglada que interseca a cada generatriz en un punto se llama curva generatriz.
		\end{defi}	
		
	\end{frame}
	
	\begin{frame}{Introducción}
		\begin{defi}
			Una superficie reglada es cilíndrica si es de la forma 
			
			\begin{equation}
			r(u,v) = \rho(u) + v a_0
			\end{equation}
			
			con $a_0 \in \R^3$.
		\end{defi}
		
		\begin{defi}
			Una superficie reglada es cónica si es de la forma 
			
			\begin{equation}
			r(u,v) = \rho_0 + v a(u)
			\end{equation}
			
			con $\rho_0 \in \R^3$.\\
			
			$\rho_0$ es el vértice del cono.
		\end{defi}
	\end{frame}
	
	\begin{frame}{Introducción}
			\begin{defi}
				Una superficie reglada es tangente desarrollable si es de la forma 
				
				\begin{equation}
				r(u,v) = \rho(u) + v \rho'(u)
				\end{equation}
			\end{defi}
			
			\begin{defi}
				Una superficie reglada que cumple que $a(u) \neq 0 \ \ \forall u \in I$ se denomina no cilíndrica. 
			\end{defi}
			
			\begin{defi}
				Una superficie no cilíndrica cuyas generatrices son paralelas a un plano directriz fijo se llama superficie de Catalan.
			\end{defi}
			
	\end{frame}
	
	\begin{frame}{Introducción}
		\begin{teo}[Caracterización de las superficies de Catalan]
			Una superficie reglada $r(u,v) = \rho(u) + v a(u)$ es una superficie de Catalan si y sólo si
			\begin{equation}
			[a(u), a'(a), a''(u)] = 0  \ \ \forall u \in I
			\end{equation}.
		\end{teo}
		
		\begin{defi}
			Una superfie de Catalan se dice conoide si todas las generatrices intersecan una recta constante (el eje del conoide).
		\end{defi}
		
	
	\end{frame}
	
	\section{Algunas superficies regladas}
	
	\begin{frame}
		\begin{center}
			\Huge\textbf{\textsf{\textcolor{blue}{Algunas \\ superficies regladas}}}
		\end{center}
	\end{frame}
	
	\subsection{Superficie cilíndrica}
	
	\begin{frame}{Superficie cilíndrica}
		
	\end{frame}
	
	\subsection{Superficie cónica}
	
	\begin{frame}{Superficie cónica}
		
	\end{frame}
	
	\subsection{Conoides}
	
	\subsubsection{Conoide de Wallis}
	
	\begin{frame}{Conoide de Wallis}
		
	\end{frame}
	
	\subsubsection{Conoide de Plücker}
	
	\begin{frame}{Conoide de Plücker}
		
	\end{frame}
	
	\subsubsection{Conoide de Plücker generalizado}
	
	\begin{frame}{Conoide de Plücker generalizado}
		
	\end{frame}
	
	\subsubsection{Paraguas de Whitney}
	
	\begin{frame}{Paraguas de Whitney}
		
	\end{frame}
	
	\subsection{Otras}
	
	\subsubsection{Banda de Möbius}
	
	\begin{frame}{Banda de Möbius}
		
	\end{frame}
	
	\subsubsection{Helicoide}
	
	\begin{frame}{Helicoide}
		
	\end{frame}
	
	\subsubsection{Hiperboloide}
	
	\begin{frame}{Hiperboloide}
		
	\end{frame}
	
	\section{Aplicaciones}
	
		\begin{frame}
			\begin{center}
				\Huge\textbf{\textsf{\textcolor{blue}{Aplicaciones}}}
			\end{center}
		\end{frame}
	
	\subsection{A la arquitectura}
	
	\subsection{Otras}
	
	
	
	
\end{document}